\section{Discussion}
\label{sec:discussion}

\FloatBarrier
\subsection{What Do the Networks Learn?}

\begin{figure}[htbp]
  \centering
  \includegraphics[width=\columnwidth]{fig5_saliency.png}
  \caption{Saliency map for \speck at 5 rounds.
  Bits 14--15 (MSB of $C_L \oplus C'_L$ difference) and bits 28--30 (corresponding position in $C_R \oplus C'_R$) carry the highest gradient magnitude, confirming that the network focuses on truncated differential patterns in the XOR-difference channel.}
  \label{fig:saliency}
\end{figure}

\begin{figure*}[htbp]
  \centering
  \includegraphics[width=\textwidth]{fig6_attention.png}
  \caption{Attention interpretability for \speck (E19).
  \emph{Left:} Bit importance heatmap across rounds---dark cells indicate high attention weight.
  \emph{Center:} Attention-augmented MLP matches baseline accuracy.
  \emph{Right:} Attention entropy peaks at 5 rounds (diffuse feature recruitment) and collapses at 7--8 rounds (noise).}
  \label{fig:attention}
\end{figure*}

Saliency analysis on \speck at 5 rounds (E08; \Cref{fig:saliency}) identifies bits 14, 15, 30, 23, and 1 as the five most important input features.
Bits 14--15 correspond to the MSB of the XOR-difference channel ($C_L \oplus C'_L$); bits 28--30 lie in the analogous position of the second ciphertext's difference channel.
This concentration on high-order difference bits is consistent with the network detecting truncated differential patterns in the penultimate round, as Benamira et al.~\citep{benamira2021deeper} hypothesized for Gohr's ResNet.

Attention-based interpretability (E19; \Cref{fig:attention}) reveals a striking entropy pattern across rounds.
At 3 rounds, attention entropy is 7.1 bits (moderately focused); it rises to 9.0 bits at 5 rounds (diffuse attention over many bits), then collapses to 6.5 bits at 7 rounds.
The interpretation: at easy rounds, the signal is concentrated in a few bits; at moderate rounds, the network recruits a broader set of features; at near-random rounds (7r, accuracy $\approx$51\%), it attends broadly but learns nothing---the low entropy at 8 rounds reflects the network giving up rather than finding signal.

\FloatBarrier
\subsection{Why Anti-Transfer?}

The anti-transfer phenomenon---accuracy below 50\% when evaluating a high-round model on low-round data---demands mechanistic explanation.

\paragraph{SPN (positive transfer).}
Gohr et al.~\citep{gohr2022assessment} proved that under independent round keys, SPN distinguishers can only learn features captured by the DDT.
DDT-based features are hierarchical: the output distribution after $r$ rounds is a convolution of the single-round DDTs, so patterns valid at $r$ rounds remain valid (though weaker) at $r' < r$.
Our \present results---99.1\% transfer from 4r to 2r---are fully consistent with this theory.

\paragraph{ARX (anti-transfer).}
For ARX ciphers, Gohr et al.~\citep{gohr2022assessment} showed that distinguishers learn non-differential features beyond the DDT.
These non-differential features exploit the interaction between modular addition and XOR, which produces round-specific correlations.
A model trained at 5 rounds learns features that \emph{anti-correlate} with the 3-round distribution: the decision boundary for 5-round ``real'' pairs actively rejects 3-round ``real'' pairs, producing the observed 45.5\% accuracy.

\paragraph{Feistel (anti-transfer---unexpected).}
Gohr et al.'s theory predicts that Feistel distinguishers, like SPN ones, should learn only differential features.
Yet we observe anti-transfer: 48.6\% for \simon 6r$\to$4r ($p < 10^{-5}$).
We hypothesize that \simon's nonlinearity---bitwise AND between rotated copies of the state---creates differential patterns that, while technically within the DDT, are highly round-dependent due to the AND operation's asymmetric propagation.
Unlike PRESENT's S-box, which achieves full diffusion within one round, \simon's Feistel structure requires multiple rounds for the AND-rotation patterns to stabilize, producing round-specific features that do not compose.
Testing this hypothesis requires training on \simon with masked intermediate rounds, which we leave to future work.

\FloatBarrier
\subsection{SIMON Key Recovery Failure}

\simon's 100\% distinguisher accuracy at 5 rounds but 0\% key recovery is the most practically informative negative result in this study.
Gohr's byte-by-byte Bayesian search assumes that partially decrypting with the correct subkey produces a distribution that the distinguisher recognizes, while wrong keys produce random-looking distributions.
In \speck, the last round's modular addition creates a direct relationship: subkey errors propagate predictably.
In \simon, the subkey enters via XOR \emph{after} the nonlinear function $f(x) = (x \lll 1) \mathbin{\&} (x \lll 8) \oplus (x \lll 2)$, meaning subkey errors interact with the state through a complex nonlinear path.
Prior work achieves \simon key recovery through different means: Tian and Hu~\citep{tian2022simon} use prepended differentials with neutral bits, bypassing the byte-by-byte search entirely.

\FloatBarrier
\subsection{Limitations}

We acknowledge five limitations:

\begin{enumerate}[leftmargin=*,itemsep=1pt]
  \item \textbf{Architecture ceiling.} Our MLP achieves fewer rounds than published SOTA (7r \speck vs.\ 8r Gohr; 9r \simon vs.\ 11r DBitNet; 6r \present vs.\ 9r DBitNet). The ResNet comparison (\Cref{tab:resnet}) shows the gap is $\leq$2.6\pp, but DBitNet and SENet may differ more substantially.
  \item \textbf{Single input difference.} We use one $\Delta P$ per cipher. Bellini et al.~\citep{bellini2023dbitnet} show that evolutionary input difference search can improve accuracy by 5+\pp.
  \item \textbf{Single-pair setting.} We use $n=2$ (one ciphertext pair). Multi-pair aggregation ($n > 2$) consistently improves accuracy in the literature~\citep{zhang2025multipair}.
  \item \textbf{Classical baseline.} Our best-bit-bias is a single-bit distinguisher. A DDT-based multi-bit distinguisher~\citep{gohr2022assessment} would reduce the reported neural--classical gap.
  \item \textbf{No PRESENT key recovery.} \present's 64-bit block requires a different partial decryption approach that we do not implement.
\end{enumerate}

\noindent
None of these limitations affect the transfer polarity finding, which is the paper's central contribution: transfer direction depends on the source and target data distributions, not on the model's absolute accuracy.
